\chapter{Hash Tables y Diccionarios}
\label{chap:hashtables}
{\huge$\bigstar$} Principiante

\section{Creación y Operaciones Básicas}
\begin{lispcode}
(defvar *ht* (make-hash-table :test 'equal))
(setf (gethash "clave" *ht*) "valor")
(gethash "clave" *ht*)
(remhash "clave" *ht*)
(clrhash *ht*)
\end{lispcode}

\section{Tests de Comparación}
\paragraph{:test} determine cómo se comparan las claves: \texttt{eq},
\texttt{eql}, \texttt{equal}, \texttt{equalp}.

\section{Iteración}
\begin{lispcode}
(maphash (lambda (k v) (format t "~A → ~A~%" k v)) *ht*)
(hash-table-count *ht*)
(hash-table-p *ht*)
\end{lispcode}

\begin{examplebox}[title=Contador de palabras]
\begin{lispcode}
(defun word-count (words)
  (let ((ht (make-hash-table :test 'equal)))
    (dolist (w words ht)
      (incf (gethash w ht 0)))))
(word-count '("the" "quick" "brown" "fox" "the"))
\end{lispcode}
\end{examplebox}

\begin{examplebox}[title=Memoización de Fibonacci]
\begin{lispcode}
(defvar *fib-cache* (make-hash-table :test 'eql))
(defun fib-fast (n)
  (or (gethash n *fib-cache*)
      (setf (gethash n *fib-cache*)
            (if (< n 2) n (+ (fib-fast (- n 1)) (fib-fast (- n 2)))))))
(fib-fast 100)
\end{lispcode}
\end{examplebox}

\begin{examplebox}[title=Two-sum con hash table]
\begin{lispcode}
(defun two-sum (nums target)
  (let ((seen (make-hash-table :test 'eql)))
    (dotimes (i (length nums))
      (let ((complement (- target (aref nums i))))
        (multiple-value-bind (found idx) (gethash complement seen)
          (when found (return-from two-sum (list idx i)))))
      (setf (gethash (aref nums i) seen) i))
    nil))
(two-sum #(2 7 11 15) 9)
\end{lispcode}
\end{examplebox}

